\ChapterImageStar[cap:marcoConceptual]{Marco Conceptual}{./images/fondo.png}\label{cap:marcoConceptual}
\mbox{}\\

\noindent
Los conceptos a continuación no solo delimitan el ámbito de estudio, sino que también proporcionan las bases terminológicas y estructurales necesarias para la evaluación, comparación e implementación de la notación propuesta para Infraestructura de \TI (\NDITI).

\noindent
\section{Notación de Modelado}
\noindent
Según~\cite{Fettke01} la notación de modelado constituye el aspecto visual y simbólico de un lenguaje de modelado, definiendo los símbolos que se utilizan para representar las construcciones abstractas de un sistema. Los autores, además distinguen este tipo de sintaxis concretas, formadas por símbolos visibles, de la sintaxis abstracta pertenecientes a los metamodelos, la cual se centra en las relaciones conceptuales entre los elementos. En el contexto de la infraestructura de TI, una notación adecuada debe permitir representar tanto la capa tecnológica (físico/virtual), como la capa lógica (servicios, datos, interconexión) y sus dependencias mutuas.

\section{Infraestructura de TI}
\noindent
La Infraestructura de~\TI~se puede definir como el conjunto de todos los elementos tecnológicos —hardware, software,  redes e instalaciones físicas— que son indispensables para el ciclo de vida de las aplicaciones y los servicios de TI, en donde es importante notar que se excluye explícitamente a las personas, los procesos de trabajo y la documentación relacionada con su gestión.~\cite{Axelos01}. Sin embargo otros autores han presentado defniciones de este concepto orientadas a una perspectiva de "Plataforma de Servicios", la cual sostiene que la infraestructura tambien abarca un conjunto de servicios organizacionales esenciales orientados a su planificación, gestión y uso por parte del factor humano\cite{Laudon01}. En contextos académicos e institucionales, el modelado de esta infraestructura facilita la comprensión, la docencia y la investigación sobre arquitecturas tecnológicas, su evolución, interoperabilidad y gobernanza.

\section{Metamodelo Conceptual}
\noindent
Un metamodelo concpetual define los conceptos, relaciones, reglas y restricciones que gobiernan la notación. por lo que pueden ser considerados la base estructural necesaria para la creación de modelos concretos de infraestructura. Como señalan~\cite{Sokolowski01}~, para las infraestructuras críticas se puede aplicar un esquema de modelado conceptual que parte de un modelo “potencial” que incluye todas las variables y relaciones genéricas que sean posibles, para luego adaptarse a casos específicos.

\section{ArchiMate}
ArchiMate es un lenguaje de modelado abierto e independiente, estandarizado por The Open Group, cuya finalidad es proporcionar una representación uniforme para la Arquitectura Empresarial (AE) ofreciendo un enfoque arquitectónico integrado que no solo describe y visualiza los diferentes dominios de la arquitectura, sino también sus relaciones y dependencias subyacentes~\cite{Vicente01}. Una de las metas centrales de este lenguaje es lograr ofrecer un enfoque integrado que describa los diferentes dominios de arquitecturay sus relaciones, lo anterior logrado a traves de una estructura compuesta de tres capas principales (Negocio, Aplicación y Tecnología)\cite{}. 
Sin embargo Archimate presenta ciertas limitaciones en su uso que ya han sido documentadas, mas especificamente~\cite{Nilus01}~en su publicacion argumenta que el lenguaje presenta ambigüedad semántica que puede originar modelos inconsistentes o engañosos; una falta de soporte general para el modelado de datos e implementaciones variadas entre las herramientas que utilizan este estándar y que en consecuencia afectan de gran manera la interoperabilidad del mismo.

\section{interoperabilidad}
\noindent
La interoperabilidad se define como la capacidad de los sistemas de usar e intercambiar datos de forma segura entre si a traves de estándares, protocolos, tecnologías y/o mecanismos para que la Información fluya~\cite{AWS01}. Estudios como los de~\cite{Bouzidi01} destacan la necesidad de definir ontologías y metamodelos comunes que permitan vincular diferentes notaciones o lenguajes sin que esto suponga la pérdida de la informacion entre modelos. En entornos académicos y de investigación, esta problemática adquiere relevancia particular, pues la falta de uniformidad en la representación de arquitecturas puede limitar la colaboración interdisciplinaria y la transferencia de conocimiento.